\subsection{A complete proof of the acyclic-support excess bound}

\begin{theorem}[Proved: eight excess arrows with acyclic support]
\label{CP2:A-high-acyclic-eight}
Let $C$ be a finite category with $k$ objects, $n$ arrows including
identities, and acyclic support.  If $n\leq2k+7$, then
$BC\not\simeq\mathbb{CP}^2$.
\end{theorem}
\begin{proof}
Suppose otherwise.  The category is connected.  With one object
its monoid has order at most nine, excluded by
Corollary~\ref{A:monoids-nine}; with two objects use
Corollary~\ref{CP2:A-high-two-objects}.  Assume $k\geq3$.
Let $E$ be the number of comparable ordered pairs of distinct
objects in the support poset.  Define
\[
 D=\sum_x(|\operatorname{End}(x)|-1),\qquad
 P=\sum_{x<y}(|C(x,y)|-1),\qquad R=E-k+1.
\]
These are nonnegative integers and
$D+P+R=n-(2k-1)\leq8$.  The support functor is onto on
fundamental groups: represent a support edge loop by
comparable edges and lift each edge to any arrow, reversing
its path when necessary.  Thus the support order complex
has zero first homology over $\mathbb F_2$.

We first record how subpieces will be used.  A clique of the
comparability graph is a totally ordered set, so its full
subcategory contains every simplex with those distinct
vertices.  A graph decomposition into clique blocks therefore
gives an actual union of nerve subcomplexes.  Across an
intersection with $H^3(-;\mathbb Q)=0$, restriction on fourth
cohomology to the two pieces is injective.  Across a rationally
acyclic connected intersection, second cohomology splits and
bounded homology passes to the pieces.  If the intersection
has zero homology in degrees at least three, boundedness
still passes in sufficiently high degrees, which is all the
Tor argument needs.  These statements follow directly from
Mayer--Vietoris, and apply successively to the clique-tree
decompositions of Lemma~\ref{CP2:A-cycle-budget}.

\paragraph{Singleton hom-sets and semisimple corners.}
If $P=0$, Lemma~\ref{CP2:A-singleton-support} pulls every
second class back from the support poset.  Its fourth
cohomology is zero by Lemma~\ref{CP2:A-high-flag}, so every
square is zero.

Suppose all corner algebras are semisimple.  If $R\leq5$,
there is no five-object chain, and the relative-bar height
bound gives $H_4=0$.  If $R\geq6$, then $D+P\leq2$.
Since $P\geq1$, either $D=0$ and $P\leq2$, or $D=1$
and $P=1$.  There is at most one five-object chain and
none longer, by Lemma~\ref{CP2:A-high-flag}.  In that
chain choose a vertex which is neither an endpoint of
a parallel consecutive hom-set nor the one nontrivial
corner, if present.  At most four vertices are forbidden
when $D=0$, and at most three when $D=1$.
The component of the degree-four relative-bar boundary
deleting this vertex is injective.  It retains all varying
arrow labels.  A composed pair of singleton arrows maps
its single basis vector to one basis vector, even if the
target hom-set is larger.  A retained nontrivial corner
has semisimple algebra, so its balanced tensor products
preserve this injection.  Endpoint deletion is the same
argument with a singleton augmentation factor.  No other
five-object chain contributes to that face component.
The fourth boundary is injective, and $H_4=0$.

We may now assume $D\geq2$ and a nonsemisimple corner.
If $P=1$, then $R\leq5$ and there is at most one corner
$M$ with nonzero $H^2$, because each such corner costs
at least four units of $D$.  If there is none, use the
parallel cut.  A singleton hom-set incident with $M$
makes its inclusion null-homotopic inside the corresponding
join, so all global second classes restrict to zero there;
the same cut again applies.  Otherwise $M$ is a leaf on
the sole doubled hom-set.  Its other endpoint has order
at most four, hence is rationally acyclic.  The deletion
of $M$ has singleton cross hom-sets and zero second
cohomology.  The two-arrow edge collage has zero products
by Lemma~\ref{CP2:A-high-two-cross}, and fourth restriction
to it and the deletion is injective.  This closes $P=1$.
Henceforth $D,P\geq2$, and $R\leq4$.

\paragraph{Tree support: $R=0$.}
There is no chain of three distinct objects: its transitive
edge would create a cycle.  The nerve is a tree of
two-object edge collages meeting at corner spaces.
If every corner has order at most five, each intersection
has zero third rational homology by
Corollary~\ref{CP2:A-five-types}.  Every edge collage has
at most
\[
 (D_x+1)+(D_y+1)+(P_{xy}+1)\leq D+P+3\leq11
\]
arrows.  Corollary~\ref{CP2:A-high-two-objects} kills its
degree-two products.  Injectivity of fourth restriction
over the tree kills the global square.

Otherwise there is a unique corner of order six or seven,
and every other corner has order at most two.  This follows
from $D+P\leq8$, $P\geq2$: a corner of order six costs
five, and a further order-three corner would cost two more.
Group all edges incident with the large corner into its
star.  The other pieces have semisimple corners and only
one strict crossing, so have no homology in degree at
least two.  They attach at rationally acyclic corners.
The star has all its arrows oriented the same way.
Collapse its small leaf corners by exact orbit functors.
Identifying the resulting trivial leaf points gives an
ordinary act cone over the large corner.  At the level
of spaces this identifies finitely many points of a
connected space, adding only a wedge of circles; the
quotient map is an isomorphism on cohomology in degrees
at least two and respects products.  The order-seven
act lemma kills its products.  The attachments cannot
restore the fourth square.

\paragraph{One cycle: $R=1$.}
The support consists of one triangle and attached trees.
If all corners have order at most five, fourth restriction
to the triangle and its bridge collages is injective.
Each bridge has at most ten arrows and zero products.
If all corners have order at most four, the intersections
are rationally acyclic, so the triangle inherits bounded
homology and second-cohomology dimension at most one.

For $D\leq3$, at most one corner has order at least three.
When it is an endpoint of the triangle, cut that rationally
acyclic corner from the opposite semisimple two-object
block, whose second cohomology vanishes.  When it is in
the middle, use Lemma~\ref{CP2:A-high-linear-Tor}.
A triangle with only semisimple corners has no fourth
homology by its height.  These arguments also cover a
triangle not containing the larger corner.

For $D=4$ or $5$ with all corners of order at most four,
$P\leq3$ or $2$, respectively.  If an adjacent hom-set
is singleton, the corresponding two-object block is a
join and has zero second cohomology; cut it from the
remaining rationally acyclic corner.  Otherwise both
adjacent hom-sets have at least two elements.  The only
possible hom-size forms, allowing smaller excess, are
$(2,2,1)$, $(2,2,2)$, $(3,2,1)$ and $(2,3,1)$, with
the long hom-set listed last.  Some extreme has both
incident stalks of size at most two.  The rank-one
small-corner lemma applies to that extreme and kills
the triangle square.

If there is a corner of order five, all other corners
have order at most two and $P\leq3$.  With the order-five
corner in the middle, collapse the two semisimple extremes.
A singleton adjacent hom-set deletes the middle object:
the relevant comma category consists of its unique arrow
to, or from, a trivial extreme.  Otherwise the preceding
four hom-size forms have a trivial extreme with all its
stalks at most two.  The objectwise two-state lemma kills
all products.  This uses no triangle-compression theorem.

If the order-five corner is an endpoint and has zero
second cohomology, cut it from the opposite semisimple
block.  If it has nonzero second cohomology, it is $P$
of Lemma~\ref{CP2:A-high-P5-graph}.  A singleton incident
hom-set is a fixed arrow, and the opposite semisimple
block has zero second cohomology; apply
Corollary~\ref{CP2:A-fixed-arrow} or its dual.
Without a singleton, both incident hom-sets have size
at least two.  Unless the long hom-set has size three,
the opposite semisimple extreme has both its stalks at
most two; collapse it and use the categorical two-state
lemma.  The sole exceptional profile is $(2,1,3)$ or
its opposite.  It consumes $P=3$ and forces both other
corners to be trivial.  Delete the middle toward the
other trivial extreme, leaving a three-state act cone
over $P$, covered by the two-row act lemma.

The remaining possibility has one corner of order six,
$D=5$, $P=2$, and all other corners trivial.  Group all
simplices containing that corner into one full piece $X$:
take its incident bridges and the triangle if it lies
there.  All remaining pieces have trivial corners and
dimension at most two, and attach to $X$ and each other
at trivial points.  Thus fourth restriction is injective.
If the corner is off the triangle, $X$ is a star and
reduces, as above, to an act cone of order six.
If it is the middle of the triangle, there is no further
neighbor, since one would create another support cycle.
A singleton adjacent hom-set deletes it; otherwise the
profile is $(2,2,1)$ and a trivial extreme has binary
stalks.  If it is an extreme, its other neighbors are on
the same side.  A singleton opposite adjacent hom-set
allows deletion of the trivial middle, leaving a star.
If that opposite hom-set is parallel, the budget forces
one of the two hom-sets incident with the order-six corner
to be singleton.  Its complement in $X$ has trivial
corners and height at most one.  The fixed-arrow lemma
therefore kills the square.  This completes $R=1$.

\paragraph{At least two cycles and $P=2$.}
Here $D\leq4$.  If all corners have order at most four,
Lemma~\ref{CP2:A-high-two-parallel} applies.  Otherwise
the unique nontrivial corner has order five, $D=4$,
and $R=2$.  If its second cohomology vanishes, the proof
of that lemma still applies: all joins have zero second
cohomology, and in the final $(2,2,1)$ triangle a trivial
extreme replaces the small-corner step.  The same argument
works if global second classes restrict to zero at the
corner, in particular if a singleton hom-set is incident
with it.  A singleton edge intersection in this argument
has a trivial endpoint and is contractible, so the fourth
restriction step remains valid.

Otherwise the corner is $P$ and every incident hom-set
is parallel.  It has one or two neighbors.  Let $X$
consist of this corner and all its neighbors, and let
$D'$ delete only the corner.  These nerves cover $NC$.
Their intersection has at most two trivial objects and
zero third homology.  The category $D'$ has trivial
corners and cycle rank at most two, hence dimension at
most two.  With one neighbor, $X$ is a two- or three-state
act cone over $P$.  With two incomparable neighbors,
it is a star and identifying the trivial leaves reduces
to a four-state act cone.  With comparable neighbors,
it is a triangle with all hom-sets of size at most two
and a trivial extreme.  Every case has zero products by
the act lemmas.  Fourth restriction finishes this case.

\paragraph{Two cycles and at least three parallel excesses.}
Now $R=2$, $D\leq3$, and $P\leq4$.  For $D=2$ there
is one order-three corner and all other corners are
trivial.  Its trivial modules on both sides have projective
dimension at most one.  To see this without classification,
use its minimal ideal: if there is one row or one column,
the corresponding trivial module is projective; the sole
nonprojective case on the other side is the two-element
one-sided band with identity, whose resolution in
Lemma~\ref{CP2:A-band-spheres} has length one on that side.
Thus every edge intersection has zero homology in degrees
at least three.  Fourth restriction to triangle pieces is
injective and they inherit bounded high homology.  At an
endpoint, the order-three corner contributes at most one
internal bar degree, so the triangle has total homological
degree at most three.  In the middle apply the small-Tor
lemma.  Consequently every triangle has zero fourth
homology.  Bridges do also, and the case is closed.

For $D=3$, all corners are rationally acyclic and at most
one has order at least three; call it the larger corner.
Single triangles and bridges are handled by the preceding
one-cycle $D\leq3$ arguments, with inherited boundedness
across corner intersections.  It remains to treat two
triangles sharing an edge, the diamond.  Write its shared
edge as $u<v$ and its incomparable exclusive vertices as
$a,b$.  Transitivity allows, up to opposition, only
$u<a,b<v$ or $a,b<u<v$.

In the first orientation, cut either endpoint from the
opposite two-edge star.  If one star has zero second
cohomology, the cut proves product vanishing.  Otherwise
each star has an edge with potentially nonzero second
cohomology.  An edge can have such cohomology only if
it is parallel and touches the larger corner: singleton
edges are joins, and semisimple edges have height one.
Thus the larger corner must be one of $a,b$, say $a$.
The shared-edge endpoints are semisimple, so that edge
has zero homology in degrees at least two.  Fourth
restriction to the two triangles is injective and both
inherit bounded high homology.  The triangle through $a$
has middle order at most four and semisimple extremes,
so has zero fourth homology by the Tor lemma.  The other
triangle is semisimple and has height two.

In the second orientation, cut $\{a,b\}$ from the shared
edge.  If the shared edge has zero second cohomology,
the cut applies.  If one exclusive vertex has singleton
hom-sets to both $u,v$, its full triangle is a join and
has zero second cohomology.  Every global second class
then restricts to zero on the shared edge, and the same
cut applies.  Avoiding these alternatives requires the
shared edge and an incident hom-set at each exclusive
vertex to be parallel.  Since $P\leq3$, these three sets
are doubled and all others are singleton.  The shared
edge touches the larger corner.  If it is $u$, collapse
the semisimple final corner $v$; all its stalks are binary,
and the objectwise two-state lemma applies.  If it is
$v$, the remaining height-one source block is semisimple
and has zero second cohomology.  The corner at $v$ has
zero second cohomology also, so the two-block lemma
finishes.  This proves every two-cycle case.

\paragraph{Three or four cycles.}
The only remaining three-cycle budget is $D=2$, $P=3$.
There is one order-three corner and all others are trivial.
The small-cycle graph decomposition consists of triangle
pieces, or one four-clique, with bridge attachments.
All edge intersections have zero homology in degree at
least three, as just proved.  Triangle pieces inherit
boundedness and have zero fourth homology by the same
endpoint-degree or middle-Tor arguments.

Consider a four-clique.  If all its hom-sets have size at
most two, at least one extreme is trivial, so the whole
piece is an objectwise binary act cone.  One hom-set of
size four is covered by the parallel cut.  The only
other possibility is one hom-set of size three and one
of size two.  Nonserial sets use the same cut.  For
separated serial sets use the downset of the second
source, exactly as in the proof of
Lemma~\ref{CP2:A-high-two-parallel}.  For adjacent serial
sets their common vertex has an additional neighbor in
the four-chain; the join-downset argument in that proof
applies.  These arguments do not depend on the sizes
two versus three, only on which hom-sets are singleton.
Thus the four-clique has zero degree-two products, and
fourth restriction kills the global square.
Finally $R=4$ forces $D=P=2$, already covered by the
two-parallel lemma.  All nonnegative budgets of sum at
most eight have now been treated.  Each contradicts the
nonzero rational square of $\mathbb{CP}^2$.
\end{proof}

\begin{corollary}[Proved: the excess-six and excess-seven assertions]
\label{CP2:A-high-six-seven}
The preceding exclusion also holds under $n\leq2k+5$ or
$n\leq2k+6$.  In particular, all acyclic-support categories
with at least seven objects and at most twenty-one arrows
are excluded.
\end{corollary}
\begin{proof}
Each displayed bound implies $n\leq2k+7$; for the last
assertion, $k\geq7$ gives $21\leq2k+7$.
\end{proof}
